% !TEX root = ../matan.tex

\section{Непрерывный путь в многомерном пространстве. Длина}

\subsection*{Пространство $\RR^d$ и его метрика}

Рассмотрим
\[
\RR^d=\bigl\{x=(x_1,\dots,x_d)\st x_k\in\RR\bigr\},\qquad d\in\NN,
\]
с евклидовым расстоянием между точками $x=(x_1,\dots,x_d)$ и
$y=(y_1,\dots,y_d)$:
\[
|x-y|:=\sqrt{\sum_{k=1}^d (x_k-y_k)^2}.
\]

\begin{Statement}{Евклидово расстояние --- метрика}{}
	Функция $(x,y)\mapsto|x-y|$ удовлетворяет аксиомам метрики:
	неотрицательность, симметрия, тождество различимости и неравенство
	треугольника.
\end{Statement}

\begin{proof}
	\emph{Неотрицательность}: $|x-y|\ge0$ как арифметический корень.
	\emph{Симметрия}: $(y_k-x_k)^2=(x_k-y_k)^2$, поэтому $|y-x|=|x-y|$.
	\emph{Различимость}: $|x-y|=0\iff\sum(x_k-y_k)^2=0\iff x_k=y_k\ \forall k\iff x=y$.
	\emph{Неравенство треугольника}: при $a_k:=x_k-y_k$, $b_k:=y_k-z_k$ имеем
	$a_k+b_k=x_k-z_k$, и по неравенству Минковского для сумм с $p=2$
	\[
	\Bigl(\sum_k (x_k-z_k)^2\Bigr)^{\!\frac12}
	\le\Bigl(\sum_k (x_k-y_k)^2\Bigr)^{\!\frac12}
	+\Bigl(\sum_k (y_k-z_k)^2\Bigr)^{\!\frac12},
	\]
	то есть $|x-z|\le|x-y|+|y-z|$.
\end{proof}

\subsection*{Непрерывный путь (кривая) и его параметризация}

\begin{Definition}{Кривая в $\RR^d$ (непрерывный путь)}{}
	\textbf{Кривой} (непрерывным путём) в $\RR^d$ называется множество точек
	\[
	\Gamma=\bigl\{\vphi(t)=(\vphi_1(t),\dots,\vphi_d(t))\bigr\}_{t\in I},
	\qquad I\subset\RR\ \text{--- промежуток},
	\]
	где все координатные функции $\vphi_k:I\to\RR$ \textbf{непрерывны}. Набор
	функций $\vphi=(\vphi_1,\dots,\vphi_d)$ задаёт \textbf{параметризацию} кривой.
\end{Definition}

\begin{Definition}{Кратная точка, простая и замкнутая кривая}{}
	Точка $x\in\RR^d$ называется \textbf{кратной}, если она отвечает по крайней
	мере двум различным значениям параметра:
	\[
	\exists\, t_1\ne t_2\in I:\quad x=\vphi(t_1)=\vphi(t_2).
	\]
	Кривая без кратных точек (концы не учитываются) называется
	\textbf{простой}; если совпадают только концы --- \textbf{замкнутой}; если
	кратных точек конечное число --- \textbf{параметризуемой}.
\end{Definition}

\subsection*{Ломаные и длина}

Длину кривой определяют как предел длин вписанных в неё ломаных.

\begin{Definition}{Ломаная}{}
	\textbf{Ломаная} --- кривая, для которой отрезок параметризации
	$I=[t_0;t_N]$ разбит на конечное число отрезков $[t_j;t_{j+1}]$ так, что на
	каждом из них все координатные функции $\vphi_k(t)$ линейны по $t$.
\end{Definition}

\begin{Definition}{Вписанная ломаная и её длина}{}
	Пусть $\Gamma=\vphi(I)$ --- кривая, а $\tau=\{t_0,t_1,\dots,t_N\}$ ---
	разбиение отрезка $I$. Ломаная $\Gamma'$ с вершинами $\vphi(t_0),\dots,\vphi(t_N)$
	(соседние вершины соединены отрезками прямой) называется \textbf{вписанной}
	в $\Gamma$. Её длина равна сумме длин звеньев:
	\[
	l(\Gamma')
	=\sum_{j=0}^{N-1}\bigl|\vphi(t_{j+1})-\vphi(t_j)\bigr|
	=\sum_{j=0}^{N-1}
	\sqrt{\sum_{k=1}^d\bigl(\vphi_k(t_{j+1})-\vphi_k(t_j)\bigr)^2}.
	\]
\end{Definition}

\begin{Definition}{Спрямляемая кривая и её длина}{}
	Кривая $\Gamma$ называется \textbf{спрямляемой}, если множество длин всех
	вписанных в неё ломаных ограничено сверху:
	\[
	\sup\bigl\{\,l(\Gamma')\st \Gamma'\ \text{--- ломаная, вписанная в }\Gamma\,\bigr\}<+\infty.
	\]
	В этом случае \textbf{длиной} кривой называется этот супремум:
	\[
	l(\Gamma):=\sup\bigl\{\,l(\Gamma')\st \Gamma'\ \text{вписана в }\Gamma\,\bigr\}.
	\]
\end{Definition}

\begin{Remark}{}{}
	Существенно, что ломаная именно \textbf{вписана} в кривую (её вершины лежат
	на $\Gamma$). При измельчении разбиения длина вписанной ломаной не убывает
	(добавление вершины не уменьшает сумму длин звеньев в силу неравенства
	треугольника), поэтому супремум можно понимать как предел по измельчающимся
	разбиениям.
\end{Remark}

\begin{Example}{Неспрямляемая кривая}{}
	Кривая
	\[
	\vphi(t)=\left(t,\ \sin\tfrac1t\right),\qquad t\in I=(0,1],
	\]
	непрерывна, но \textbf{не} спрямляема. 
	Действительно, возьмём точки
	\[
	t_n=\frac{1}{\tfrac{\pi}{2}+\pi n},\qquad n=0,1,2,\dots,
	\]
	в которых $\sin\frac1{t_n}=\sin\!\bigl(\tfrac{\pi}{2}+\pi n\bigr)=(-1)^n$.
	Для соответствующей вписанной ломаной с вершинами $\vphi(t_n)$ длина каждого
	звена не меньше модуля приращения второй координаты:
	\[
	\bigl|\vphi(t_n)-\vphi(t_{n+1})\bigr|\ge\bigl|(-1)^n-(-1)^{n+1}\bigr|=2.
	\]
	Поэтому ломаная, построенная по первым $N+1$ из этих точек, имеет длину
	$\ge 2N$, что неограниченно растёт при $N\to\infty$. Значит, супремум длин
	вписанных ломаных равен $+\infty$, и кривая не спрямляема.
\end{Example}

